% =============================================================================
% Section 7: Statistical Analysis
% =============================================================================

\section{Statistical Analysis}
\label{sec:statistical}

% AGENT: Fill from systematics fitter Wave 4 report
% Document the statistical framework, likelihood construction, and fit methodology.

\subsection{Likelihood Model}
\label{sec:statistical:likelihood}

% AGENT: Fill likelihood model description from WorkspaceBuilder output.
% Include number of channels, bins, samples, and nuisance parameters.

The statistical analysis is based on a binned profile likelihood ratio test.
The likelihood function is constructed as:
\begin{equation}
  \mathcal{L}(\mu, \boldsymbol{\theta}) =
  \prod_{c \in \text{channels}} \prod_{b \in \text{bins}}
  \text{Pois}\!\left(n_{cb} \;\middle|\; \mu s_{cb}(\boldsymbol{\theta})
  + \sum_p b_{pcb}(\boldsymbol{\theta})\right)
  \cdot \prod_k p(\tilde{\theta}_k | \theta_k)
\end{equation}
where $\mu$ is the signal strength parameter, $\boldsymbol{\theta}$ are the
nuisance parameters, and $p(\tilde{\theta}_k | \theta_k)$ are the constraint
terms.

% AGENT: Fill workspace structure details from pyhf workspace JSON.

The \texttt{pyhf} HistFactory workspace contains:
\begin{itemize}
  \item Channels: [list channels]
  \item Signal samples: [count]
  \item Background samples: [count]
  \item Systematic sources: [count] (after pruning)
  \item Total bins: [count]
\end{itemize}

\subsection{Fit Methodology}
\label{sec:statistical:fit}

% AGENT: Fill fit configuration from Fitter settings.
% Include test statistic, CL method, POI range, and convergence criteria.

The test statistic $\tilde{q}_\mu$ is used for upper limit calculations,
defined as:
\begin{equation}
  \tilde{q}_\mu = -2 \ln \frac{\mathcal{L}(\mu, \hat{\hat{\boldsymbol{\theta}}})}
  {\mathcal{L}(\hat{\mu}, \hat{\boldsymbol{\theta}})}
  \qquad \text{with } 0 \leq \hat{\mu} \leq \mu
\end{equation}

Upper limits are computed using the CL$_\text{s}$ method with asymptotic
approximation. The signal strength parameter $\mu$ is scanned in the range
$[0, \mu_\text{max}]$.

\subsection{Expected Sensitivity}
\label{sec:statistical:sensitivity}

% AGENT: Fill expected sensitivity from SensitivityOptimizer results.
% Include expected limit value, median and +/-1,2 sigma bands.

The expected sensitivity is evaluated using an Asimov dataset (background-only
hypothesis). The expected 95\% CL upper limit on the signal strength is:
\begin{equation}
  \mu_\text{exp}^{95} = [value] \quad
  \left( [+1\sigma], [-1\sigma] \right)
\end{equation}

% AGENT: Fill interpretation in terms of cross-section if applicable.
